\section{System Model}\label{sec:model}
\subsection{Multiple-Access Channel}
We consider the sensing architecture in Fig.~\ref{fig:system}. Let $K\in\mathbb N$ denote the number of transmitters and let $I=[-1,1]\subset\R$. Transmitter $k\in\{1,\ldots,K\}$ holds a calibrated real observation $s_k\in I$. We write $\mathbf s=(s_1,\ldots,s_K)^\T\in I^K\subset\R^K$. All transmitters apply the same continuous encoder $\enc:I\to\C$ and transmit simultaneously over one complex channel use. We identify $\C$ with $\R^2$: the real and imaginary components are the in-phase and quadrature coordinates of the transmitted symbol. The inputs and the desired function remain real valued. Unless stated otherwise, our converse and optimal-design results assume $K\ge4$.

For a given $P>0$, we impose the per-transmitter pointwise power constraint
\begin{equation}
 x_k=\enc(s_k),\qquad \sup_{s\in I}|\enc(s)|^2\le P.
 \label{eq:peak-model}
\end{equation}
Thus every admissible observation produces a symbol satisfying the transmit-power limit. We assume synchronized transmissions and unit effective channel gains after compensation. The receiver observes
\begin{equation}
 Y=\sum_{k=1}^{K}x_k+N=r_\enc(\mathbf s)+N,
 \qquad \widehat f(Y)=D(Y),
 \label{eq:mac-model}
\end{equation}
where we define the noiseless superposition $r_\enc:I^K\to\C$ by $r_\enc(\mathbf s)=\sum_{k=1}^K\enc(s_k)$ and use a receiver $D:\C\to\R$. The disturbance $N\in\C$ is arbitrary in our uniform error analysis. For the small-noise analysis, we specify $N\sim\mathcal{CN}(0,N_0)$, independent of $\mathbf s$, with $\mathbb E|N|^2=N_0$. The receiver estimates the desired function directly from the superposition. Calibration and synchronization are assumed resources; their overhead is excluded from the one-use payload count. Section~\ref{sec:implementation} quantifies bounded residual implementation and channel errors.

\subsection{Desired Functions and Examples}
We fix the input calibration and desired function before designing the encoder. For continuous real-valued functions $\phi,\psi:I\to\R$, we consider $f_{\phi,\psi}:I^K\to\R$ defined by
\begin{equation}
 f_{\phi,\psi}(\mathbf s)
 =\frac1K\sum_{k=1}^K\phi(s_k)
   -\psi\!\left(\frac1K\sum_{k=1}^Ks_k\right),
 \qquad \mathbf s\in I^K.
 \label{eq:target}
\end{equation}
We use $\phi$ as a local preprocessing function and $\psi$ as a function of the empirical mean. Define the real sums $S=\sum_k s_k$ and $T=\sum_k\phi(s_k)$, and their normalized statistics $m=S/K\in I$ and $t=T/K\in\R$. Thus $f_{\phi,\psi}(\mathbf s)=t-\psi(m)$. For integer $p\ge1$, let $t_p=K^{-1}\sum_k s_k^p$ be the empirical $p$th raw moment. In particular, $Q=\sum_k s_k^2$ and $q=Q/K=t_2$.

\emph{Mean.} Choosing $\phi(s)=s$ and $\psi\equiv0$ gives $f(\mathbf s)=m$. The encoder $\enc(s)=\sqrt P\,s$ and noiseless receiver $D(y)=\operatorname{Re}(y)/(K\sqrt P)$ compute the mean using one real additive statistic. This example belongs to \eqref{eq:target}, but does not satisfy the nonaffineness and strict-convexity assumptions of our converse. We do not claim that the mean requires two transmitted statistics.

\emph{Variance.} Choosing $\phi(s)=\psi(s)=s^2$ gives the empirical variance
$$
 V(\mathbf s)=q-m^2=\frac1K\sum_{k=1}^K(s_k-m)^2.
$$
For example, $\enc(s)=\sqrt{P/2}\,(s+i s^2)$ obeys \eqref{eq:peak-model}. In the absence of noise, $Y_0=r_\enc(\mathbf s)$ yields
$$
 m=\frac{\operatorname{Re}Y_0}{K\sqrt{P/2}},\qquad
 q=\frac{\operatorname{Im}Y_0}{K\sqrt{P/2}},\qquad V=q-m^2.
$$
This construction explains how two real moments can be conveyed in one complex channel use. Section~\ref{sec:examples} optimizes their scaling and the translation of the transmitted curve.

\emph{Jensen gaps.} If $\psi=\phi$ is convex, \eqref{eq:target} is the empirical Jensen gap. The choices $\phi(s)=s^4$ and $\phi(s)=e^s$ give $t_4-m^4$ and $K^{-1}\sum_k e^{s_k}-e^m$, respectively. The quantity $t_4-m^4$ is not the fourth central moment, which also involves the second and third raw moments. More generally, the boundedness of $\phi$ makes $\enc(s)=\alpha[s+i\phi(s)]$ feasible for $\alpha=\sqrt{P/(1+\|\phi\|_\infty^2)}>0$, so both $S$ and $T$ are sufficient for exact computation. Our converse asks when every exact shared continuous encoder must preserve both sums.

\emph{Relation to nomographic functions.} A scalar nomographic representation has the form
$$
 f(\mathbf s)=\Phi\!\left(\sum_{k=1}^K h_k(s_k)\right),
$$
where $h_k:I\to\R$ are local preprocessing functions and $\Phi$ is a scalar postprocessing function defined on the attainable sums~\cite{goldenbaum2015}. In contrast, our construction uses two real additive statistics, $(S,T)\in\R^2$, followed by $F(m,t)=t-\psi(m)$ with $(m,t)=(S/K,T/K)$. It is therefore a two-coordinate version of nomographic computation, implemented through the two coordinates of one complex symbol. These classes overlap: the mean is already scalar nomographic. Under the assumptions of Theorem~\ref{thm:rigidity}, however, no shared continuous scalar encoder $h:I\to\R$ can compute \eqref{eq:target} exactly, because the theorem forces two linearly independent real coefficient vectors. This conclusion concerns the shared continuous-encoder model; it does not exclude representations with different preprocessing functions at different transmitters or discontinuous preprocessing.

\subsection{Exact Computation and Separation}
For a specified continuous desired function $f:I^K\to\R$, we call $\enc$ \emph{exact} if
\begin{equation}
 \forall\,\mathbf s,\mathbf t\in I^K,\qquad
 r_\enc(\mathbf s)=r_\enc(\mathbf t)
 \ \Longrightarrow\ f(\mathbf s)=f(\mathbf t).
 \label{eq:exactness}
\end{equation}
Equivalently, there is a well-defined noiseless receiver $D_0:\mathcal R_\enc\to\R$, where $\mathcal R_\enc=r_\enc(I^K)\subset\C$, such that $D_0(r_\enc(\mathbf s))=f(\mathbf s)$ for every $\mathbf s\in I^K$. This receiver is continuous: the continuous surjection from the compact set $I^K$ onto the Hausdorff set $\mathcal R_\enc$ is a quotient map. Exactness specifies which function values the superposition must distinguish; the distances between these superpositions determine noise sensitivity.

For a nonconstant $f:I^K\to\R$, we define the squared-separation score
\begin{equation}
 \Lambda_{\sq}(\enc;f)
 =\inf_{\substack{\mathbf s,\mathbf t\in I^K\\f(\mathbf s)\ne f(\mathbf t)}}
 \frac{|r_\enc(\mathbf s)-r_\enc(\mathbf t)|^2}
 {|f(\mathbf s)-f(\mathbf t)|^2}\in[0,\infty).
 \label{eq:score}
\end{equation}
A larger score means that unequal values of the desired function produce better-separated received sums. In particular, $\Lambda_{\sq}>0$ guarantees $|f(\mathbf s)-f(\mathbf t)|\le |r_\enc(\mathbf s)-r_\enc(\mathbf t)|/\sqrt{\Lambda_{\sq}}$ for every $\mathbf s,\mathbf t\in I^K$ and implies exactness. Proposition~\ref{prop:amplification} identifies $1/\sqrt{\Lambda_{\sq}}$ as the smallest uniform noise-amplification constant achievable by a receiver.

We maximize \eqref{eq:score} over the class
$$
 \mathcal E_P=\{\enc\in C(I;\C):\sup_{s\in I}|\enc(s)|^2\le P\},\qquad
 \Lambda_{\sq}^\star(f)=\sup_{\enc\in\mathcal E_P}\Lambda_{\sq}(\enc;f).
$$
Under the assumptions below, Section~\ref{sec:design} proves that every exact encoder has the form $\enc(s)=c+us+v\phi(s)$, giving $r_\enc(\mathbf s)=Kc+uS+vT$. The offset $c$ places the curve inside the peak-power disk; the lengths and relative angle of $u,v$ control the sensitivity of moment recovery. This result reduces the optimization over continuous functions to a finite-dimensional choice of encoder coefficients.
